\documentclass[11pt]{article}
\usepackage[margin=1in]{geometry}
\usepackage[T1]{fontenc}
\usepackage{lmodern}
\usepackage{microtype}
\usepackage{amsmath}
\usepackage{booktabs}
\usepackage{tikz}
\usetikzlibrary{arrows.meta,positioning,calc}
\usepackage[numbers]{natbib}
\setlength{\bibsep}{3pt}
\usepackage[hidelinks]{hyperref}

\newcommand{\code}[1]{\texttt{#1}}

\title{Assignment 1: A Pentium M style hybrid branch predictor\\
  \large COSC 6385 Computer Architecture, Fall 2026}
\author{Bryan E. Tuck\\University of Houston\\{\small \url{https://github.com/Bryan-Tuck/cosc6385-hw1-branch-predictor}}}
\date{October 2026}

\begin{document}
\maketitle

\section{Introduction}

The Pentium M predicts branch direction with a bimodal table, a global predictor and a loop
predictor \cite{uzelac2009}. This assignment covers the first two. The bimodal table of 2-bit
counters, indexed by the branch address, holds each branch's usual direction. The global
predictor, a set-associative table of tagged 2-bit counters, holds directions that depend on
the path taken to the branch. The Pentium M indexes it with the branch address XORed with a
path information register (PIR). The assignment does not require the PIR, so my predictor,
like its Appendix~B example, XORs the address with a global history register (GHR) of
conditional branch outcomes instead. When a set holds an entry whose tag matches, the global counter is
used; otherwise the bimodal counter is used and the global predictor allocates an entry for
the next encounter.

I implemented this hybrid as the \code{pm\_predictor} class of \code{my\_predictor.h} in the CBP2
branch prediction infrastructure \cite{jimenez2006}, following the structure and the worked
example in the assignment \cite{assignment}, and measured it on the 20 CBP2 traces. It averages
7.759 mispredictions per thousand instructions (MPKI) over the traces (Section~\ref{sec:results}).
The loop predictor and branch target prediction are outside the assignment and are not implemented,
and I did not attempt the extra credit \code{cpm\_predictor}. The predictor uses no third-party
cache code.

\section{Predictor design}

Figure~\ref{fig:design} shows the structure the code implements. The sizes follow the assignment
figure: a 4096-entry bimodal table, a 4-way global predictor with 512 sets, and
$\mathrm{Index} = \mathrm{HASH}[14{:}6]$, $\mathrm{Tag} = \mathrm{HASH}[5{:}0]$. Each size is a
\code{\#define} at the top of the \code{pm\_predictor} section of \code{my\_predictor.h}.

\begin{figure}[t]
\centering
\begin{tikzpicture}[
  >=Latex, font=\small,
  box/.style={draw, rounded corners=1pt, minimum height=2.2em, inner sep=4pt, align=center},
  tbl/.style={draw, minimum width=3.8em, minimum height=4em, align=center},
]
  \node[box] (pc) {branch address $A$};
  \node[box, below=1.4em of pc] (ghr) {GHR $H$ (15 bits)};
  \node[draw, circle, inner sep=1pt, right=1.8em of ghr] (xor) {$\oplus$};
  \node[box, right=1.4em of xor] (split) {$h = A \oplus H$\\$h[14{:}6]$ = set\\$h[5{:}0]$ = tag};
  \node[tbl, right=1.4em of split] (w0) {way 0};
  \node[tbl, right=0.3em of w0] (w1) {way 1};
  \node[tbl, right=0.3em of w1] (w2) {way 2};
  \node[tbl, right=0.3em of w2] (w3) {way 3};
  \node[above=0.3em of w1.north east] {global predictor, 512 sets, entry = V, tag, 2bc, LRU age};
  \node[box, below=1.4em of w1.south east] (hit) {valid and tag match in some way?};
  \node[box, below=1.2em of hit] (mux) {hit ? global counter : bimodal counter};
  \node[tbl, minimum width=8em] (bim) at (split |- mux) {bimodal table\\4096 $\times$ 2bc};
  \draw[->] (pc.east) -| (xor.north);
  \draw[->] (ghr) -- (xor);
  \draw[->] (xor) -- (split);
  \draw[->] (split) -- (w0) node[midway, above, font=\scriptsize] {set};
  \draw[->] (pc.west) -- ++(-1.0em,0) |- (bim.west) node[pos=0.75, above, font=\scriptsize] {$A[11{:}0]$};
  \draw[->] (w1.south east) -- (hit.north) node[midway, right, font=\scriptsize] {4 tags, 4 counters};
  \draw[->] (hit) -- (mux);
  \draw[->] (bim.east) -- (mux.west);
  \draw[->] (mux.east) -- ++(1.4em,0) node[right] {prediction};
\end{tikzpicture}
\caption{Structure of \code{pm\_predictor}. The hash $h$ of the branch address and the global
history register selects a set of the 4-way global predictor and supplies the tag; the bimodal
table is indexed by the low address bits. A tag hit selects the global prediction, a miss
selects the bimodal prediction.}
\label{fig:design}
\end{figure}

\subsection{Bimodal table}

The bimodal table holds $2^{12} = 4096$ saturating 2-bit counters. For a conditional branch at
address $A$ the counter index is the low 12 bits,
%
\begin{equation}
  i_{\mathrm{bim}} = A \bmod 2^{12},
\end{equation}
%
and the prediction is the counter's upper bit (states 2 and 3 predict taken). The counters
start at 0. The index is $\mathrm{IP}[11{:}0]$, as in the assignment figure. The Appendix~B
example drops the two low bits of its word-aligned addresses, but x86 branch instructions can
start at any byte, so the index keeps bits [1:0].

\subsection{Global predictor}

The global predictor has four ways of 512 sets. Each entry, a \code{pm\_global\_entry} in the
code, stores a valid bit, a 6-bit tag, a 2-bit counter and a 2-bit LRU age. The global history
register $H$ holds the outcomes of the last 15 conditional branches, newest in bit 0. The lookup
hashes the address with the history and splits the hash into a set index and a tag:
%
\begin{align}
  h &= (A \oplus H) \bmod 2^{15}, \\
  s &= \lfloor h / 2^{6} \rfloor = h[14{:}6], \qquad t = h \bmod 2^{6} = h[5{:}0].
\end{align}
%
The upper hash bits select the set and the lower bits form the tag, following the labels
$\mathrm{Index} = \mathrm{HASH}[14{:}6]$ and $\mathrm{Tag} = \mathrm{HASH}[5{:}0]$ in the
assignment figure and the row/tag split of its Appendix~B example, where a 4-bit hash gives the
row from bits [3:2] and the tag from bits [1:0]. Because the 15-bit history covers the whole
hash, the tag depends on the six most recent outcomes XORed with address bits [5:0], and the
set index on the older nine outcomes XORed with address bits [14:6].

The lookup compares $t$ with the tag of each valid way in set $s$. A match is a hit and the
global prediction is the upper bit of that way's counter. On a miss the lookup records the LRU way of the set
(age $= 3$) as the victim for allocation. Replacement is true LRU kept with per-entry ages. When
a way is used, every way of the set whose age was smaller than the used way's age grows by one
and the used way's age becomes 0. In the constructor, way $w$ gets age $3 - w$, so an empty set
fills way 0, 1, 2, 3 in that order and no valid entry is evicted while an invalid way remains.
For two ways the age is a single bit that is 1 in the LRU way, which is the LRU column of the
Appendix~B tables.

\subsection{Selection and update}

The final prediction is the global prediction on a hit and the bimodal prediction on a miss.
There is no chooser table; the tag comparison selects the table. The predictor returns a
\code{pm\_update} object that records the bimodal index, the set, the tag, the hit flag and the
way, so \code{update} does not repeat the lookup.

After the driver reports the outcome, \code{update} changes state only for conditional branches:
\begin{enumerate}
  \item The bimodal counter at $i_{\mathrm{bim}}$ is trained toward the outcome whether the
    global lookup hit or missed.
  \item On a hit the matching global counter is trained toward the outcome. On a miss the LRU
    way is allocated. Its valid bit is set, its tag becomes $t$, and its existing counter is
    trained toward the outcome without being reset, as in the Appendix~B example, where in
    step 9 the evicted entry's counter 3 becomes 2 after a not-taken outcome.
  \item The used way becomes most recently used, as described above.
  \item The outcome is shifted into bit 0 of $H$, which is masked to 15 bits.
\end{enumerate}
Unconditional branches are predicted taken and do not touch the tables or the history, the same
convention as the sample \code{gshare\_predictor} that ships with the infrastructure. Branch
targets are not predicted; \code{predict} sets the target prediction to 0, as the starter code
does.

\section{Results}
\label{sec:results}

Each CBP2 trace holds the branches of 100 million instructions of one SPEC CPU2000 or SPECjvm98
program. The driver counts direction mispredictions of conditional branches and reports MPKI,
mispredictions per thousand instructions; lower is better. Table~\ref{tab:mpki} lists the
result for every trace. The raw numbers are in \code{results/hybrid.csv}.

\begin{table}[htbp]
\centering
\caption{MPKI of \code{pm\_predictor} per trace (lower is better).}
\label{tab:mpki}
\begin{tabular}{lr@{\hspace{3em}}lr}
\toprule
Trace & MPKI & Trace & MPKI \\
\midrule
164.gzip & 12.697 & 213.javac & 3.001 \\
175.vpr & 15.969 & 222.mpegaudio & 2.932 \\
176.gcc & 18.009 & 227.mtrt & 3.453 \\
181.mcf & 18.269 & 228.jack & 4.239 \\
186.crafty & 9.109 & 252.eon & 2.153 \\
197.parser & 10.936 & 253.perlbmk & 5.095 \\
201.compress & 8.019 & 254.gap & 4.309 \\
202.jess & 2.005 & 255.vortex & 1.643 \\
205.raytrace & 3.602 & 256.bzip2 & 0.115 \\
209.db & 4.801 & 300.twolf & 24.821 \\
\midrule
\multicolumn{3}{l}{Average over the 20 traces} & 7.759 \\
\bottomrule
\end{tabular}
\end{table}

MPKI ranges from 0.115 on bzip2 to 24.821 on twolf. The four highest, twolf, mcf, gcc and vpr,
are all SPECint programs.

\bibliographystyle{plainnat}
\bibliography{refs}

\end{document}
